\chapter{Häufige Fehler}
\label{ch:quality/common-mistakes}

Dieses Kapitel beschreibt häufige Fehler und wie du sie vermeidest.

\section{Vorsicht beim Speichern von
ObjectRefs}\label{vorsicht-beim-speichern-von-objectrefs}

Ein ObjectRef\index{objectref@\texttt{ObjectRef}} (also ein Spieler oder
eine Entity)\index{entity@Entity} wird ungültig, wenn der Spieler oder
die Entity, für die es steht, das Spiel verlässt. Das passiert, wenn der
Spieler offline geht oder die Entity entladen oder entfernt wird.

Seit Luanti 5.2 geben die Methoden eines ungültigen ObjectRefs immer
\texttt{nil} zurück; jeder Aufruf wird im Grunde ignoriert.

Speichere ObjectRefs möglichst nicht. Wenn du es doch tust, prüfe vor
der Verwendung, ob das ObjectRef noch gültig ist:

\begin{Shaded}
\begin{Highlighting}[]
\CommentTok{{-}{-} Funktioniert erst ab Luanti 5.2}
\ControlFlowTok{if} \VariableTok{obj}\OperatorTok{:}\NormalTok{get\_pos}\OperatorTok{()} \ControlFlowTok{then}
    \CommentTok{{-}{-} ist gültig!}
\ControlFlowTok{end}
\end{Highlighting}
\end{Shaded}

\section{Formspec-Eingaben nicht
vertrauen}\label{formspec-eingaben-nicht-vertrauen}

Bösartige Clients können Formspecs\index{formspec@Formspec} jederzeit
und mit beliebigem Inhalt abschicken.

Der folgende Code hat zum Beispiel eine Sicherheitslücke, durch die
Spieler sich selbst Moderatorrechte geben können:

\begin{Shaded}
\begin{Highlighting}[]
\KeywordTok{local} \KeywordTok{function}\NormalTok{ show\_formspec}\OperatorTok{(}\VariableTok{name}\OperatorTok{)}
    \ControlFlowTok{if}
        \KeywordTok{not} \VariableTok{core}\OperatorTok{.}\NormalTok{check\_player\_privs}\OperatorTok{(}\VariableTok{name}\OperatorTok{,} \OperatorTok{\{} \VariableTok{privs} \OperatorTok{=} \KeywordTok{true} \OperatorTok{\})}
    \ControlFlowTok{then}
        \ControlFlowTok{return} \KeywordTok{false}
    \ControlFlowTok{end}

    \VariableTok{core}\OperatorTok{.}\NormalTok{show\_formspec}\OperatorTok{(}
        \VariableTok{name}\OperatorTok{,}
        \StringTok{"modman:modman"}\OperatorTok{,}
        \VerbatimStringTok{[[}
\VerbatimStringTok{        size[3,2]}
\VerbatimStringTok{        field[0,0;3,1;target;Name;]}
\VerbatimStringTok{        button\_exit[0,1;3,1;sub;Promote]}
\VerbatimStringTok{    ]]}
    \OperatorTok{)}
    \ControlFlowTok{return} \KeywordTok{true}
\KeywordTok{end}

\VariableTok{core}\OperatorTok{.}\NormalTok{register\_on\_player\_receive\_fields}\OperatorTok{(}
    \KeywordTok{function}\OperatorTok{(}\VariableTok{player}\OperatorTok{,} \VariableTok{formname}\OperatorTok{,} \VariableTok{fields}\OperatorTok{)}
        \CommentTok{{-}{-} SCHLECHT! Hier fehlt die Prüfung der Privilegien!}

        \KeywordTok{local} \VariableTok{privs} \OperatorTok{=} \VariableTok{core}\OperatorTok{.}\NormalTok{get\_player\_privs}\OperatorTok{(}\VariableTok{fields}\OperatorTok{.}\VariableTok{target}\OperatorTok{)}
        \VariableTok{privs}\OperatorTok{.}\VariableTok{kick} \OperatorTok{=} \KeywordTok{true}
        \VariableTok{privs}\OperatorTok{.}\VariableTok{ban} \OperatorTok{=} \KeywordTok{true}
        \VariableTok{core}\OperatorTok{.}\NormalTok{set\_player\_privs}\OperatorTok{(}\VariableTok{fields}\OperatorTok{.}\VariableTok{target}\OperatorTok{,} \VariableTok{privs}\OperatorTok{)}
        \ControlFlowTok{return} \KeywordTok{true}
    \KeywordTok{end}
\OperatorTok{)}
\end{Highlighting}
\end{Shaded}

Mit einer Prüfung der Privilegien\index{privileg@Privileg} ist die Lücke
geschlossen:

\begin{Shaded}
\begin{Highlighting}[]
\VariableTok{core}\OperatorTok{.}\NormalTok{register\_on\_player\_receive\_fields}\OperatorTok{(}
    \KeywordTok{function}\OperatorTok{(}\VariableTok{player}\OperatorTok{,} \VariableTok{formname}\OperatorTok{,} \VariableTok{fields}\OperatorTok{)}
        \KeywordTok{local} \VariableTok{name} \OperatorTok{=} \VariableTok{player}\OperatorTok{:}\NormalTok{get\_player\_name}\OperatorTok{()}
        \ControlFlowTok{if}
            \KeywordTok{not} \VariableTok{core}\OperatorTok{.}\NormalTok{check\_player\_privs}\OperatorTok{(}
                \VariableTok{name}\OperatorTok{,}
                \OperatorTok{\{} \VariableTok{privs} \OperatorTok{=} \KeywordTok{true} \OperatorTok{\}}
            \OperatorTok{)}
        \ControlFlowTok{then}
            \ControlFlowTok{return} \KeywordTok{false}
        \ControlFlowTok{end}

        \CommentTok{{-}{-} Code}
    \KeywordTok{end}
\OperatorTok{)}
\end{Highlighting}
\end{Shaded}

\section{ItemStacks nach dem Ändern
zurücksetzen}\label{itemstacks-nach-dem-uxe4ndern-zuruxfccksetzen}

Ist dir aufgefallen, dass es in der API einfach
\texttt{ItemStack}\index{itemstack@\texttt{ItemStack}} heißt und nicht
\texttt{ItemStackRef}, analog zu \texttt{InvRef}? Ein \texttt{ItemStack}
ist nämlich keine Referenz, sondern eine Kopie. Du arbeitest mit einer
Kopie der Daten, nicht mit dem Stack im
Inventar.\index{inventar@Inventar} Änderst du den Stack, ändert sich der
Stack im Inventar also nicht.

Mach das zum Beispiel nicht:

\begin{Shaded}
\begin{Highlighting}[]
\KeywordTok{local} \VariableTok{inv} \OperatorTok{=} \VariableTok{player}\OperatorTok{:}\NormalTok{get\_inventory}\OperatorTok{()}
\KeywordTok{local} \VariableTok{stack} \OperatorTok{=} \VariableTok{inv}\OperatorTok{:}\NormalTok{get\_stack}\OperatorTok{(}\StringTok{"main"}\OperatorTok{,} \DecValTok{1}\OperatorTok{)}
\VariableTok{stack}\OperatorTok{:}\NormalTok{get\_meta}\OperatorTok{():}\NormalTok{set\_string}\OperatorTok{(}\StringTok{"description"}\OperatorTok{,} \StringTok{"Angebissen"}\OperatorTok{)}
\CommentTok{{-}{-} SCHLECHT! Die Änderung geht verloren}
\end{Highlighting}
\end{Shaded}

Sondern so:

\begin{Shaded}
\begin{Highlighting}[]
\KeywordTok{local} \VariableTok{inv} \OperatorTok{=} \VariableTok{player}\OperatorTok{:}\NormalTok{get\_inventory}\OperatorTok{()}
\KeywordTok{local} \VariableTok{stack} \OperatorTok{=} \VariableTok{inv}\OperatorTok{:}\NormalTok{get\_stack}\OperatorTok{(}\StringTok{"main"}\OperatorTok{,} \DecValTok{1}\OperatorTok{)}
\VariableTok{stack}\OperatorTok{:}\NormalTok{get\_meta}\OperatorTok{():}\NormalTok{set\_string}\OperatorTok{(}\StringTok{"description"}\OperatorTok{,} \StringTok{"Angebissen"}\OperatorTok{)}
\VariableTok{inv}\OperatorTok{:}\NormalTok{set\_stack}\OperatorTok{(}\StringTok{"main"}\OperatorTok{,} \DecValTok{1}\OperatorTok{,} \VariableTok{stack}\OperatorTok{)}
\CommentTok{{-}{-} Richtig! Der ItemStack wird gesetzt}
\end{Highlighting}
\end{Shaded}

Bei Callbacks\index{callback@Callback} ist es etwas komplizierter.
Änderst du einen übergebenen \texttt{ItemStack}, ändert er sich auch für
den Aufrufer und alle folgenden Callbacks. In der Engine gespeichert
wird er aber nur, wenn der Aufrufer des Callbacks ihn setzt.

\begin{Shaded}
\begin{Highlighting}[]
\VariableTok{core}\OperatorTok{.}\NormalTok{register\_on\_item\_eat}\OperatorTok{(}
    \KeywordTok{function}\OperatorTok{(}
        \VariableTok{hp\_change}\OperatorTok{,}
        \VariableTok{replace\_with\_item}\OperatorTok{,}
        \VariableTok{itemstack}\OperatorTok{,}
        \VariableTok{user}\OperatorTok{,}
        \VariableTok{pointed\_thing}
    \OperatorTok{)}
        \VariableTok{itemstack}
            \OperatorTok{:}\NormalTok{get\_meta}\OperatorTok{()}
            \OperatorTok{:}\NormalTok{set\_string}\OperatorTok{(}\StringTok{"description"}\OperatorTok{,} \StringTok{"Angebissen"}\OperatorTok{)}
        \CommentTok{{-}{-} Fast richtig! Die Daten gehen verloren,}
        \CommentTok{{-}{-} wenn ein anderer Callback das Verhalten abbricht}
    \KeywordTok{end}
\OperatorTok{)}
\end{Highlighting}
\end{Shaded}

Bricht kein Callback ab, wird der Stack gesetzt und die Beschreibung
aktualisiert. Bricht aber ein Callback ab, kann die Änderung verloren
gehen.

Besser ist es so:

\begin{Shaded}
\begin{Highlighting}[]
\VariableTok{core}\OperatorTok{.}\NormalTok{register\_on\_item\_eat}\OperatorTok{(}
    \KeywordTok{function}\OperatorTok{(}
        \VariableTok{hp\_change}\OperatorTok{,}
        \VariableTok{replace\_with\_item}\OperatorTok{,}
        \VariableTok{itemstack}\OperatorTok{,}
        \VariableTok{user}\OperatorTok{,}
        \VariableTok{pointed\_thing}
    \OperatorTok{)}
        \VariableTok{itemstack}
            \OperatorTok{:}\NormalTok{get\_meta}\OperatorTok{()}
            \OperatorTok{:}\NormalTok{set\_string}\OperatorTok{(}\StringTok{"description"}\OperatorTok{,} \StringTok{"Angebissen"}\OperatorTok{)}
        \VariableTok{user}\OperatorTok{:}\NormalTok{get\_inventory}\OperatorTok{():}\NormalTok{set\_stack}\OperatorTok{(}
            \StringTok{"main"}\OperatorTok{,}
            \VariableTok{user}\OperatorTok{:}\NormalTok{get\_wield\_index}\OperatorTok{(),}
            \VariableTok{itemstack}
        \OperatorTok{)}
        \CommentTok{{-}{-} Richtig, die Beschreibung wird immer gesetzt!}
    \KeywordTok{end}
\OperatorTok{)}
\end{Highlighting}
\end{Shaded}

Brechen die Callbacks ab oder setzt der Aufrufer den Stack nicht, wird
die Änderung trotzdem übernommen. Setzen die Callbacks oder der Aufrufer
den Stack, schadet das zusätzliche \texttt{set\_\allowbreak{}stack}
nicht.
